\section{Etapa 2 --- Investigaci\'on asistida}
\label{sec:et2}
\periodo{10 --- 15 de abril de 2026 \quad$\cdot$\quad \textasciitilde{}119 \emph{commits} en 4 d\'ias}

Etapa corta en calendario y enorme en volumen. Se defini\'o un procedimiento de
trabajo ---en la jerga de la herramienta, una \emph{skill} llamada
\texttt{INVESTIGADOR}--- y se lo dej\'o correr de forma aut\'onoma sobre el vault
de investigaci\'on durante dos d\'ias.

\begin{figure}[ht]
\centering
\includegraphics[width=0.96\textwidth]{comparacion_metodos.png}
\caption{S\'intesis visual del espacio de m\'etodos que organiz\'o la investigaci\'on
asistida. Resume funci\'on y log\'istica; no establece una jerarqu\'ia universal.}
\end{figure}

\subsection{Qu\'e se le pidi\'o al agente}

El procedimiento fijaba cuatro cosas, y vale la pena registrarlas porque son
las que hicieron que el resultado fuera utilizable y no ruido:

\begin{enumerate}[itemsep=0.2em]
\item \textbf{Naturaleza anal\'itica}: no resumir por resumir, sino explicar el
mecanismo f\'isico.
\item \textbf{Gesti\'on de fuentes}: todo lo que se afirme debe rastrearse a la
base de 65 papers ya construida en la Etapa~1 ---sin fuente, no entra.
\item \textbf{Arquitectura en Obsidian}: un concepto por archivo, con bloque de
definici\'on, ejemplo concreto tomado de un paper real, y enlaces salientes.
\item \textbf{Criterio cient\'ifico}: distinguir lo establecido de lo
especulativo.
\end{enumerate}

\subsection{Qu\'e produjo}

\begin{center}\small
\begin{tabularx}{\textwidth}{@{}X r@{}}
\toprule
\textbf{Resultado} & \textbf{Cantidad} \\
\midrule
Conceptos at\'omicos enriquecidos (definici\'on + ejemplo + enlaces) & 65 \\
Conceptos at\'omicos nuevos & \textasciitilde{}37 \\
Total de conceptos en el vault al cierre & \textasciitilde{}140 \\
Ejemplos concretos por cap\'itulo (uno por cada uno de los 8) & 4 c/u \\
\emph{Wikilinks} rotos al cerrar la verificaci\'on & 0 \\
Archivos hu\'erfanos (sin ning\'un enlace entrante) & 0 \\
\bottomrule
\end{tabularx}
\end{center}

Los conceptos nuevos se agruparon en tres bloques que anticipan exactamente lo
que despu\'es hubo que implementar:

\begin{itemize}[itemsep=0.2em]
\item \textbf{Procesamiento de se\~nales}: transformadas, ventaneo,
\emph{phase-shift} ---base del \texttt{masw\_dispersion.py} de la Etapa~7.
\item \textbf{Inversi\'on avanzada}: no unicidad, regularizaci\'on, b\'usqueda
global ---justificaci\'on del uso de un algoritmo evolutivo en la Etapa~7.
\item \textbf{Se\~nales discretas}: muestreo, \emph{aliasing}, cuantizaci\'on
---base del dise\~no del ADC y la decimaci\'on en la Etapa~3.
\end{itemize}

\subsection{Consolidaci\'on y cierre}

El 15 de abril se revis\'o manualmente lo producido, se reorganizaron las
im\'agenes en \texttt{Assets/} y se cerr\'o la rama. Este paso importa: el
resultado del agente no se acept\'o autom\'aticamente, se revis\'o antes de
integrarlo.

\seleccionado{\textbf{Qued\'o:} el vault de \textasciitilde{}140 conceptos
enlazados, que funcion\'o despu\'es como referencia r\'apida durante todo el
dise\~no ---sobre todo los conceptos de procesamiento y de inversi\'on. Tambi\'en
qued\'o el m\'etodo: definir el criterio primero, dejar correr, revisar despu\'es.}

\pendiente{El vault no se volvi\'o a auditar despu\'es de abril. Algunos
conceptos pueden haber quedado desactualizados respecto de decisiones
posteriores ---por ejemplo, todo lo que se escribi\'o suponiendo compensaci\'on
por cancelaci\'on polo-cero, enfoque que se abandon\'o (Etapa~3).}
